\documentclass[UTF8,a4paper,12pt]{ctexart}

\usepackage{geometry}
\usepackage{array}
\usepackage{tabularx}
\usepackage{longtable}
\usepackage{booktabs}
\usepackage{enumitem}
\usepackage{hyperref}
\usepackage{xcolor}
\usepackage{listings}
\usepackage{fancyhdr}

\geometry{left=2.6cm,right=2.6cm,top=2.5cm,bottom=2.5cm}
\setlength{\headheight}{15pt}
\setlist{nosep,leftmargin=2em}
\hypersetup{
  colorlinks=true,
  linkcolor=black,
  urlcolor=blue
}

\pagestyle{fancy}
\fancyhf{}
\lhead{软件测试作业二：GUI 测试}
\rhead{测试计划}
\cfoot{\thepage}

\lstset{
  basicstyle=\ttfamily\small,
  breaklines=true,
  frame=single,
  columns=fullflexible,
  backgroundcolor=\color{gray!8}
}

\begin{document}

\begin{center}
  {\LARGE\bfseries 测试计划}
\end{center}
\thispagestyle{empty}
\vspace{1em}
\tableofcontents
\newpage

\section{文档说明}

本文档为软件测试课程作业二《GUI 测试》的测试计划，用于指导 SauceDemo 电商网站的 Selenium 自动化测试设计、脚本编写、测试执行和后续测试报告整理。

\section{测试目标}

本次测试的主要目标如下：

\begin{enumerate}
  \item 使用 Selenium 等 GUI 自动化测试框架模拟用户在 Web 页面中的真实操作。
  \item 验证 SauceDemo 网站核心购物流程是否能够正常完成。
  \item 覆盖两级以上页面跳转，包括登录页、商品列表页、商品详情页、购物车页、结算页和订单完成页。
  \item 设计不少于 5 个完整功能或事务的测试场景，且用户登录不作为独立功能统计，只作为业务流程的前置步骤。
  \item 覆盖作业要求中的按钮、输入框、表单、下拉框、链接等必选 Web 元素，并至少覆盖一种可选元素。
  \item 使用大模型或软件开发 agent 辅助完成测试计划、测试脚本和测试报告，提高测试设计与实现效率。
\end{enumerate}

\section{被测对象}

被测网站为 SauceDemo 电商网站。

网站地址：

\begin{lstlisting}
https://www.saucedemo.com/
\end{lstlisting}

测试账号：

\begin{lstlisting}
username: standard_user
password: secret_sauce
\end{lstlisting}

选择 SauceDemo 作为被测对象的原因如下：

\begin{enumerate}
  \item SauceDemo 是典型的 Web 电商演示系统，包含商品浏览、商品详情、购物车、结算和订单完成等业务流程。
  \item 网站页面层级清晰，能够从登录页跳转到商品列表页、商品详情页、购物车页、结算页和订单完成页，满足“两级以上页面”的要求。
  \item 页面中包含按钮、输入框、表单、下拉框、链接、侧边栏导航等多种 Web 元素，适合开展 GUI 自动化测试。
  \item 网站不涉及真实支付、真实用户隐私和复杂验证码，适合作为课程作业的自动化测试对象。
\end{enumerate}

\section{测试范围}

本次测试围绕 SauceDemo 的核心购物流程展开：

\begin{lstlisting}
登录页 -> 商品列表页 -> 商品详情页 -> 购物车页 -> 结算页 -> 订单完成页
\end{lstlisting}

测试范围包含：

\begin{enumerate}
  \item 用户登录后进入商品浏览流程。
  \item 商品列表页商品展示与排序。
  \item 商品详情页跳转与商品信息展示。
  \item 商品加入购物车与购物车数量更新。
  \item 购物车页面商品信息检查与结算入口跳转。
  \item 结算信息填写、订单确认与订单完成。
  \item 侧边栏导航菜单打开、关闭和菜单项可见性检查。
  \item 移动端浏览器环境下登录、页面滚动和加入购物车等关键交互的轻量抽测。
\end{enumerate}

测试范围不包含：

\begin{enumerate}
  \item 真实支付流程测试。
  \item 后端接口性能测试。
  \item 数据库验证。
  \item 安全渗透测试。
  \item 多浏览器深度兼容性测试。
  \item 非测试账号相关的权限体系测试。
\end{enumerate}

\section{测试环境与工具}

\begin{table}[htbp]
\centering
\caption{测试环境与工具}
\begin{tabularx}{\textwidth}{>{\bfseries}p{3.6cm}X}
\toprule
项目 & 内容 \\
\midrule
操作系统 & Windows 10/11 \\
浏览器 & Google Chrome \\
自动化框架 & Selenium WebDriver \\
编程语言 & Python 3.x \\
测试运行框架 & pytest \\
浏览器驱动管理 & webdriver-manager 或 Selenium Manager \\
大模型/agent 工具 & Codex、Copilot 或其他软件开发 agent \\
辅助材料 & 测试计划、测试代码、测试报告、大模型使用记录截图 \\
\bottomrule
\end{tabularx}
\end{table}

建议 Python 依赖如下：

\begin{lstlisting}
selenium
pytest
webdriver-manager
\end{lstlisting}

移动端 Web UI 测试方式：

\begin{enumerate}
  \item 优先使用 Appium 驱动 Android 模拟器或真实 Android 设备中的 Chrome 浏览器，打开 SauceDemo Web 网站并执行移动端 UI 自动化测试。
  \item 移动端测试作为桌面端 Web 测试的轻量补充，不重复覆盖完整购物流程，重点关注触控点击、页面滚动、移动端布局、关键元素可见性和加入购物车状态变化。
  \item 回退方案：如果 Appium、Android 模拟器或真机环境搭建受限，则使用 Chrome mobile emulation 或 Selenium 设置移动端窗口尺寸进行响应式 Web 兼容性验证，并在测试报告中说明该方式属于移动端视口模拟，不等同于完整移动端 UI 测试。
\end{enumerate}

\section{测试策略}

本次测试采用功能测试与 GUI 自动化测试结合的策略：

\begin{enumerate}
  \item 先根据网站核心业务流程设计测试场景和测试用例。
  \item 再使用 Selenium 将测试场景转化为自动化测试脚本。
  \item 使用显式等待提升脚本稳定性，避免由于页面加载时间导致测试误判。
  \item 对关键页面跳转、关键元素状态和关键业务结果进行断言。
  \item 对可重复执行的流程使用 pytest 组织测试方法。
  \item 测试执行后记录通过、失败、异常信息和必要截图。
  \item 使用大模型辅助生成测试计划、测试脚本初稿和测试报告内容，再由小组成员检查、修改和运行验证。
\end{enumerate}

\section{测试功能与事务设计}

用户登录不作为独立功能统计，只作为以下业务事务的前置步骤。

\begingroup
\scriptsize
\setlength{\tabcolsep}{2pt}
\begin{longtable}{p{1.1cm}p{2.2cm}p{2.7cm}p{2.7cm}p{3.8cm}p{1.2cm}}
\caption{测试功能与事务设计}\\
\toprule
编号 & 功能或事务 & 页面层级/跳转 & 主要操作 & 预期结果 & 负责人 \\
\midrule
\endfirsthead
\toprule
编号 & 功能或事务 & 页面层级/跳转 & 主要操作 & 预期结果 & 负责人 \\
\midrule
\endhead
TC01 & 登录后进入商品浏览流程 & 登录页 -> 商品列表页 & 输入账号、密码，点击 Login 按钮 & 成功进入商品列表页，页面显示商品列表和购物车入口 & 成员 A \\
TC02 & 商品排序 & 商品列表页 & 使用排序下拉框选择 Name (Z to A)、Price (low to high) 等选项 & 商品列表按照所选规则重新排序 & 成员 A \\
TC03 & 查看商品详情 & 商品列表页 -> 商品详情页 & 点击商品名称或商品图片链接 & 进入商品详情页，显示商品名称、描述、价格和 Add to cart 按钮 & 成员 B \\
TC04 & 添加商品到购物车 & 商品列表页/商品详情页 -> 购物车状态变化 & 点击 Add to cart 按钮 & 按钮状态变化为 Remove，购物车 badge 数量正确增加 & 成员 B \\
TC05 & 购物车页面检查 & 商品列表页 -> 购物车页 & 点击购物车链接，查看购物车商品 & 购物车页显示已添加商品的名称、价格、数量和 Checkout 按钮 & 成员 B \\
TC06 & 结算与订单完成 & 购物车页 -> 结算信息页 -> 订单确认页 -> 完成页 & 点击 Checkout，填写结算信息，继续并完成订单 & 订单完成页显示成功提示，购物流程结束 & 成员 C \\
TC07 & 侧边栏导航菜单 & 商品列表页 -> 侧边栏显示/隐藏 & 点击菜单按钮，检查导航项，点击关闭按钮 & 侧边栏正常展开和关闭，导航项可见 & 成员 C \\
TC08 & 移动端 Web UI 轻量抽测 & Android Chrome 中登录后商品列表页 & 使用 Appium 驱动移动端 Chrome，验证登录、商品列表滑动和加入购物车状态变化 & 页面元素可见，触控点击和滚动正常，购物车 badge 可正确更新 & 成员 C \\
\bottomrule
\end{longtable}
\endgroup

以上测试功能数量不少于 5 个，并且覆盖多页面跳转与完整业务事务，满足作业要求。

\section{Web 元素覆盖情况}

\begingroup
\small
\setlength{\tabcolsep}{2pt}
\begin{longtable}{p{1.8cm}p{1.8cm}p{6.8cm}@{\hspace{0.4cm}}p{3.4cm}}
\caption{Web 元素覆盖情况}\\
\toprule
元素类型 & 作业要求 & SauceDemo 中的对应元素 & 对应用例 \\
\midrule
\endfirsthead
\toprule
元素类型 & 作业要求 & SauceDemo 中的对应元素 & 对应用例 \\
\midrule
\endhead
按钮 & 必选 & 登录按钮（Login）、加入购物车按钮（Add to cart）、移除按钮（Remove）、结算按钮（Checkout）、继续按钮（Continue）、完成按钮（Finish）、侧边栏关闭按钮 & TC01、TC04、TC05、TC06、TC07 \\
输入框 & 必选 & 用户名输入框（username）、密码输入框（password）、姓名输入框（first name、last name）、邮政编码输入框（postal code） & TC01、TC06 \\
表单 & 必选 & 登录表单、结算信息表单 & TC01、TC06 \\
下拉框 & 必选 & 商品排序下拉框 & TC02 \\
链接 & 必选 & 商品名称链接、商品图片链接、购物车链接、侧边栏菜单项 & TC03、TC05、TC07 \\
侧边栏/导航栏 & 可选 & 左侧菜单按钮、全部商品（All Items）、关于页面（About）、退出登录（Logout）、重置应用状态（Reset App State） & TC07 \\
购物车数量标记 & 可选 & 购物车图标旁显示的商品数量标记（badge） & TC04、TC05 \\
移动端触控/滑动 & 可选 & Android Chrome 中的触控点击、页面滚动与元素定位 & TC08 \\
\bottomrule
\end{longtable}
\endgroup

本次计划覆盖按钮、输入框、表单、下拉框、链接 5 类必选元素，并覆盖侧边栏/导航栏、购物车 badge、移动端触控与滑动等可选元素，满足作业对 Web 元素种类的要求。

\section{测试用例设计要点}

\subsection{TC01 登录后进入商品浏览流程}

前置条件：浏览器打开 SauceDemo 登录页。

测试步骤：

\begin{enumerate}
  \item 在 username 输入框输入 \texttt{standard\_user}。
  \item 在 password 输入框输入 \texttt{secret\_sauce}。
  \item 点击 Login 按钮。
  \item 检查是否进入商品列表页。
  \item 检查商品列表、排序下拉框和购物车入口是否可见。
\end{enumerate}

预期结果：页面跳转到商品列表页，URL 或页面标题包含 inventory 相关信息，商品列表正常显示。

\subsection{TC02 商品排序}

前置条件：用户已成功进入商品列表页。

测试步骤：

\begin{enumerate}
  \item 定位商品排序下拉框。
  \item 选择 Name (Z to A)。
  \item 获取排序后的商品名称列表。
  \item 再选择 Price (low to high)。
  \item 获取排序后的商品价格列表。
\end{enumerate}

预期结果：商品名称或价格按照所选规则正确排列。

\subsection{TC03 查看商品详情}

前置条件：用户已成功进入商品列表页。

测试步骤：

\begin{enumerate}
  \item 点击某个商品名称链接。
  \item 跳转到商品详情页。
  \item 检查商品名称、描述、价格、商品图片和 Add to cart 按钮。
  \item 点击 Back to products 返回商品列表页。
\end{enumerate}

预期结果：商品详情页信息完整，返回按钮可正常回到商品列表页。

\subsection{TC04 添加商品到购物车}

前置条件：用户已成功进入商品列表页或商品详情页。

测试步骤：

\begin{enumerate}
  \item 点击某商品的 Add to cart 按钮。
  \item 检查按钮文本是否变为 Remove。
  \item 检查购物车 badge 是否显示数量 1。
  \item 可继续添加第二个商品并检查数量变化。
\end{enumerate}

预期结果：商品成功加入购物车，购物车数量与添加商品数量一致。

\subsection{TC05 购物车页面检查}

前置条件：购物车中已有至少一个商品。

测试步骤：

\begin{enumerate}
  \item 点击购物车链接。
  \item 进入购物车页面。
  \item 检查商品名称、价格、数量。
  \item 检查 Continue Shopping 和 Checkout 按钮。
\end{enumerate}

预期结果：购物车页面显示的商品信息与已添加商品一致，结算入口可用。

\subsection{TC06 结算与订单完成}

前置条件：购物车中已有商品，并已进入购物车页。

测试步骤：

\begin{enumerate}
  \item 点击 Checkout 按钮。
  \item 在 first name、last name、postal code 输入框填写结算信息。
  \item 点击 Continue 按钮。
  \item 在订单确认页检查商品、价格和汇总信息。
  \item 点击 Finish 按钮。
\end{enumerate}

预期结果：页面进入订单完成页，并显示订单完成提示信息。

\subsection{TC07 侧边栏导航菜单}

前置条件：用户已成功进入商品列表页。

测试步骤：

\begin{enumerate}
  \item 点击页面左上角菜单按钮。
  \item 检查侧边栏是否展开。
  \item 检查 All Items、About、Logout、Reset App State 等菜单项是否可见。
  \item 点击关闭按钮。
\end{enumerate}

预期结果：侧边栏能够正常展开和关闭，导航菜单项显示正常。

\subsection{TC08 移动端 Web UI 轻量抽测}

前置条件：优先搭建 Appium 测试环境，并启动 Android 模拟器或连接真实 Android 设备；设备中的 Chrome 浏览器可访问 SauceDemo 网站。若 Appium 或移动端设备环境搭建受限，则回退为 Chrome mobile emulation 或 Selenium 移动端窗口尺寸验证，并在测试报告中说明回退原因。

测试步骤：

\begin{enumerate}
  \item 使用 Appium 启动 Android Chrome 并打开 SauceDemo 登录页。
  \item 通过触控方式输入账号和密码，点击 Login 按钮完成登录。
  \item 检查商品列表页主要元素在移动端屏幕中是否可见。
  \item 滑动商品列表页面，检查商品卡片和按钮在移动端滚动场景下是否可见。
  \item 点击一个商品的 Add to cart 按钮，检查购物车 badge 是否更新为 1。
\end{enumerate}

预期结果：移动端浏览器环境下登录页和商品列表页关键元素可见，触控输入、触控点击、页面滑动和加入购物车状态变化正常。移动端测试仅作为轻量补充抽测，不重复执行桌面端已经覆盖的商品详情、购物车检查和完整结算流程。

\end{document}
